\section{Weighted mixtures under relaxed host conditions}
\label{sec:weighted}

For an abelian $F\supset K$, define
\[
 \alpha_F(p)=
 \begin{cases}e_p^{-1},&f_p=1,\\0,&f_p>1,
 \end{cases}
 \qquad (p\equiv1\pmod3).
                                                        \tag{9.1}
\]
Lemma~\ref{lem:forced-splitting} says that its normalized
degree-one cubic trace is $a_p\alpha_F(p)$. Thus $0\le\alpha_F\le1$,
including primes ramified in $F$.

\begin{proposition}[Positive mixtures]
If $c_i\ge0$, $\sum_i c_i=1$, and
$\sum_i c_i\alpha_{F_i}(p)=1$ on a set of split-$K$ primes,
then every field with $c_i>0$ splits completely at every prime
in that set.
\end{proposition}
\begin{proof}
For each prime, $\sum_i c_i(1-\alpha_{F_i}(p))=0$ is a sum of
nonnegative numbers. Each active summand is zero, so
$\alpha_{F_i}(p)=1$, which means $e_p=f_p=1$.
The same proof applies to a countable nonnegative mixture.
\end{proof}

For literal hosts Theorem~\ref{thm:empty-norms} is stronger and
does not require positivity. The next estimate remains useful
after the common cyclotomic condition is removed.

\begin{proof}[Proof of Theorem~\ref{thm:mixture-cost}]
Let $D=[F:\Q]$ and write $\Lambda_F(n)$ for the coefficients of
$-\zeta_F'/\zeta_F$. Abelian factorization expresses $\zeta_F$
as the product of the primitive Dirichlet $L$-functions of its
$D$ characters. Apply the uniform smooth explicit formula proved
in Section~6 to each character, and divide by $D$. The single
principal character supplies the main term. The
conductor--discriminant formula gives
\begin{align}
 \frac1D\sum_n\Lambda_F(n)W(n/x)
 &=\frac{xI}{D}
       +O_W\!\left(\sqrt{x}\left(1+
                    \frac1D\sum_\xi\log\cond(\xi)\right)\right)\notag\\
 &=\frac{xI}{D}+O_W\bigl(\sqrt{x}(1+\log\rd(F))\bigr),
 \qquad I=\int_1^2W(u)\,du>0.
                                                        \tag{9.2}
\end{align}
The implied constant is independent of $D$, $F$, and the number
of characters. Complex characters cause no problem: the estimates
are absolute and their sum is the real expression on the left.

Every $\Lambda_F(n)$ is nonnegative. At a split-$K$ prime $p$,
its coefficient at $n=p$, divided by $D$, is exactly
$\alpha_F(p)\log p$: if $f_p=1$, there are $D/e_p$ prime ideals
of norm $p$; otherwise there are none. Discarding all other
terms in (9.2) yields
\[
 \sum_{p\equiv1(3)}\alpha_F(p)\log p\,W(p/x)
 \le\frac{xI}{D}+C_W\sqrt{x}(1+\log\rd(F)).
                                                        \tag{9.3}
\]
With $D\ge c x^2$ and $\log\rd(F)\le x^\epsilon/72$, this is
at most $C_{W,c}x^{1/2+\epsilon}$ for $x\ge1$.

Put $A(x)=\sum_{p\equiv1(3)}\log p\,W(p/x)$.
The prime number theorem in this fixed progression gives
$A(x)\sim Ix/2$. For an exact mixture, the triangle inequality
and (9.3) give
\[
 A(x)\le\sum_i|c_i|
       \sum_{p\equiv1(3)}\alpha_{F_i}(p)\log p\,W(p/x)
       \le C_{W,c}x^{1/2+\epsilon}\sum_i|c_i|.
                                                        \tag{9.4}
\]
This proves the claimed lower bound. If the weighted absolute
defect is at most $A(x)/4$, the first left side in (9.4) may
instead be $3A(x)/4$, and the conclusion is unchanged.
\end{proof}

This is a cost bound for pointwise recovery or weighted absolute
approximation. Applying only separate absolute estimates to the
host terms incurs this coefficient mass. It does not assert that
one particular signed aggregate cannot cancel or admit a different
comparison. The unconditional zero statement for literal hosts
has the stronger, but precisely ordinary-Euler-product, scope of
Theorem~\ref{thm:empty-norms}.

\section{A fixed quotient and its exact Artin projector}
\label{sec:projector}

Put $F_0=K\Q(\zeta_{16})^+$ and
$G_0=\Gal(F_0/K)\simeq C_4$. The factors in
Theorem~\ref{thm:projector} are finite-order Hecke $L$-functions.
Their entireness and functional equations follow from Tate's
theory; see \citet[Theorem~5.22]{Poonen2015}. No factor is principal,
because $\chi$ has order three and $\psi$ has order dividing four.

\begin{proof}[Proof of Theorem~\ref{thm:projector}]
Let $G=\Gal(F/K)$, $H=\Gal(F/F_0)$, and $r=|H|$. Since
$\chi$ is already a character of $G_K$, the tensor identity for
induction gives
\[
 V=\operatorname{Ind}_{G_F}^{G_K}(\chi|_{G_F})
        \simeq\chi\otimes\mathbb C[G].
                                                        \tag{10.1}
\]
For clarity, the isomorphism is the usual coset model: after
multiplying the component at a representative $g$ by $\chi(g)$,
the induced action becomes the scalar $\chi(\gamma)$ times
left translation by the image of $\gamma$ in $G$.
Right translations commute with this action. Averaging them over
$H$ gives the orthogonal idempotent
\[
          P_H=\frac1r\sum_{h\in H}R_h.
                                                        \tag{10.2}
\]
Its image is $\chi\otimes\mathbb C[G/H]$, of dimension four.
The regular representation of the abelian quotient is the direct
sum of its four characters. Taking the determinant of
$1-\operatorname{Frob}_{\mathfrak p}T$ on inertia invariants,
prime by prime, proves
\[
 J_0(s)=L_{F_0}(s,\chi_{F_0})
        =\prod_{\psi\in\widehat G_0}L_K(s,\chi\psi).
                                                        \tag{10.3}
\]
The same argument proves the asserted trace identity at every
prime power. At a ramified prime use $V^{I_{\mathfrak p}}$;
$P_H$ commutes with inertia, and
$(\operatorname{im}P_H)^{I_{\mathfrak p}}
=P_H(V^{I_{\mathfrak p}})$.

For the integral assertion, let $\mathcal O=\Z[\omega]$.
For each coset $C$ of $H$ in $G$, put $v_C=\sum_{g\in C}e_g$.
These four vectors are a basis of $\mathcal O[G]^H$, and
$v_C/r$ are a basis of $P_H\mathcal O[G]$. The quotient is
therefore $(\mathcal O/r\mathcal O)^4$; its determinant ideal
is $(r)^4$ and its cardinality is $r^8$. This proves the exact
index, including the coefficient-ring normalization. Each $v_C$
has Euclidean norm $\sqrt r$, and $P_H$ has norm one.

Finally $L_K(s,\chi)$ occurs as the factor $\psi=1$, and every
other factor is entire, so its zeros cannot be cancelled in
$J_0$. Thus RH for $J_0$ implies RH for the target. The dual
collection equals the original collection: conjugating ideals
of $K$ changes $\chi^{-1}$ into $\chi$, while the characters
coming from the abelian field $\Q(\zeta_{16})^+$ are unchanged
under this conjugation and are permuted by inversion. Its
functional equation therefore also shows that zero-freeness
for $\Re s>1/2$ suffices for RH.
\end{proof}

At an unramified split-$K$ rational prime the first coefficient
of $-J_0'/J_0$ is $4a_p\log p$ if $p\equiv\pm1\pmod{16}$,
and is zero otherwise. The potentially visible split classes
are exactly $1,31$ modulo $48$. Hence recovering all eight
split classes individually is stronger than necessary for a
fixed product that contains the target factor.

\subsection{An explicit sufficient analytic comparison}

This fixed product gives a precise possible replacement for an
ordinary host transfer. Let
\[
 -J_0'(s)/J_0(s)=\sum_{n\ge1}c(n)n^{-s},\qquad
 S_0(x)=\sum_{n\ge1}c(n)e^{-n/x}.
                                                        \tag{10.4}
\]
The finite-image degree-eight Artin representation gives
$|c(n)|\le8\Lambda(n)$, including at ramified prime powers.
If a construction proved $S_0(x)=O_\eta(x^{1/2+\eta})$ for every
$\eta>0$, it would prove RH for the target. Indeed, for $\Re s>1$,
absolute convergence and the substitution $u=n/x$ give
\[
 \int_0^\infty S_0(x)x^{-s-1}\,dx
                       =-\Gamma(s)J_0'(s)/J_0(s).
                                                        \tag{10.5}
\]
The proposed bound, together with exponential decay as $x\to0$,
makes the left side holomorphic on $\Re s>1/2$. Since
$\Gamma$ has no zeros or poles there, the logarithmic derivative
has no poles, and the entire function $J_0$ has no zeros there.
Its functional equation completes the assertion. The bound in
this paragraph is a sufficient condition, not an estimate proved
by the present paper or by the abstract projector.

\subsection{Normalization, denominators, and gain}

The coefficients in (10.2) have absolute sum one relative to the
\emph{unnormalized} twisted traces
$\tr(R_h\operatorname{Frob}^j\mid V)$. If these traces have
already been divided by $|G|$, recovering (10.3) instead has
coefficient mass $|G|$. Norm one is not a free normalization
of an estimate on normalized host traces.

For a general $H$-stable torsion-free $\Z$-lattice $M$ with
$a=\operatorname{rank}_{\Z}M^H$, averaging gives
\[
 M^H\subset P_HM\subset r^{-1}M^H,
 \qquad [P_HM:M^H]\le r^a.
                                                        \tag{10.6}
\]
The inclusions follow directly by multiplying the average by
$r$. Thus the denominator cost is at most $a\log r$, not
unconditionally $\log r$ for a large theta module. In the
rank-four regular example the logarithmic cost divided by
$[\Q(\omega):\Q]=2$ is exactly $4\log r$.

Orthogonal projection preserves nonnegativity of the trace of
a commuting positive operator. It does not make Frobenius
positive. A determinant bound alone does not force its divisibility
to occur on the invariant component. For an exact counterexample
take $M=\Z e_+\oplus\Z e_-$ with a $C_2$-action of signs
$+1,-1$. The equivariant integral map
\[
                 T=\operatorname{diag}(1,p^N)
                                                        \tag{10.7}
\]
has full cokernel of order $p^N$ and invariant-component
cokernel of order one, for every $N\ge1$.

Uniform lattice inclusion is different and gives a constructive
inheritance statement. If $f:M_{\mathbb Q}\to N_{\mathbb Q}$,
projectors $P,Q$ satisfy $fP=Qf$, and
$f(M)\subset\pi^a N$ over a discrete valuation ring, then
\[
                    f(PM)=Qf(M)\subset\pi^a QN.
                                                        \tag{10.8}
\]
Thus the same valuation gain survives on the projected lattices.
There is no new lower-gain theorem to prove for such an inclusion
once the coherent maps and lattices have been constructed. Comparing
$PM,QN$ with invariant integral lattices may still incur the
denominator cost in (10.6).

This distinction applies to the source. Its \texttt{vGeLat}
asserts that a vector lies in
$\pi^{\lceil g v_p(p)\rceil}N$, and \texttt{LatInst.vKge}
applies this to the image of the chain map. In a discrete valuation
ring that scaled lattice is independent of the chosen uniformizer.
A statement for every vector in a given source lattice is therefore
a uniform inclusion of the form above, rather than merely a
cokernel determinant. The unresolved task is to construct compatible
projectors on those source and target lattices, preserve the
relevant source restriction, and identify the required trace and
normalization. If the known inclusion applies on that source,
its valuation automatically passes to the projected lattices.

\section{The highest weights prevent the naive deck action}
\label{sec:highest}

We now prove the specific obstruction in
Theorem~\ref{thm:highest}. All assertions concern one fixed
family with scale-independent $C_c$. The source has
$G=\Gal(F/\Q)$ of order $4n$ and a CM-label subgroup
$\Psi\subset G$ of index two. Its top set is
$\mathcal T=\{\tau_1(t):t\in TT\}\subset\Psi$, with
$\tau_1$ injective. The source size and weight conditions give
\begin{align}
 0<x^{-\epsilon/100}n\delta
     &\le |\mathcal T|\le2n\delta,
       &\delta&=x^{-\epsilon/8},\notag\\
 a_T(t)&\ge a=x\exp(x^\epsilon),
       &S_c&=4k\Delta,\notag\\
 kS_c&\le w''(n),
       &w''(n)&=\frac{a}{xn\exp(x^{\epsilon/3})}.
                                                        \tag{11.1}
\end{align}
Here $\Delta$ is the source positive integer denoted
\texttt{Del}. At a top label its first Harish--Chandra
coordinate is $a_T(t)+(k-1)/2$. At a non-top label its first
coordinate $u_0$ and the adjacent coordinate $v_0$ satisfy
\[
     u_0-v_0\le C_cS_c,\qquad 0\le v_0\le S_c.
                                                        \tag{11.2}
\]
Every word member has the same first coordinate. These are the
equalities \texttt{top\_u}, \texttt{cloud\_gap},
\texttt{cloud\_v}, and \texttt{hcW\_first}; the highest
weight is obtained by subtracting the fixed $\rho$-vector.

\begin{proof}[Proof of Theorem~\ref{thm:highest}]
For $H\subset G$ of index at most eight, put $H_0=H\cap\Psi$.
Then $|H_0|\ge n/4$. Multiplication makes $H_0$ act freely
on $\Psi$, so every orbit has at least $n/4$ elements.
For sufficiently large $x$, $\delta<1/8$; (11.1) then says
$0<|\mathcal T|<n/4$. Thus $\mathcal T$ is not a union of
$H_0$-orbits. Some $h\in H_0$ moves a top label to a non-top label.

At a top label the first highest-weight coordinate is exactly
$a_T(t)\ge x\exp(x^\epsilon)$. At a non-top label it is
$u_0-(k-1)/2\le(C_c+1)S_c$. Put
$C_+=\max(1,C_c+1)$. Since $k,n\ge1$, the last inequality
in (11.1), which is a clause of \texttt{SizeDel}, gives
\[
 \frac{S_c}{a}\le\frac1{kxn\exp(x^{\epsilon/3})}<\frac1x.
                                                        \tag{11.3}
\]
For $x>C_+$, the first highest-weight coordinate at every top
label is strictly greater than at every non-top label. For all
sufficiently large $x$ the earlier cardinality condition also holds.
In particular the top set can be recovered from the
weight tuple itself, and that tuple is not $H_0$-invariant.

The source coefficient module, after complexification, is
\[
       W=\bigotimes_{\tau\in\Psi}V_{\lambda_\tau}
       \quad\text{for}\quad
       \mathcal G=\prod_{\tau\in\Psi}\mathrm{GL}_k(\mathbb C),
                                                        \tag{11.4}
\]
with the determinant shifts included. The irreducible modules
$V_{\lambda_\tau}$ have their indicated highest weights;
see \citet[Sections~20.3--20.8]{Milne2018}. The product
representation has a unique highest-weight line for the product
Borel. An intertwiner between it and the representation obtained
by permuting the factors by $h$ must carry this line to the
highest-weight line. Equality of their diagonal-torus characters
then forces $\lambda_{h\tau}=\lambda_\tau$ for every $\tau$.
The element found above violates this equality. Therefore no
action of all of $H_0$ can implement this coordinate permutation
on the one fixed coefficient representation.
\end{proof}

\subsection{What an equivariant enlargement would still require}

The conclusion concerns the natural field-permutation action,
not all possible linear actions. A direct sum over conjugated
weight tuples can formally carry an action permuting the
summands. This is an action relating different representations;
it does not by itself identify a commuting Galois action with
the Artin representation in (10.1), or a projected theta gain.
Also the reference complex embedding in the literal data can
change between transported objects. Although the family fixes
the labelled top set, this relabelling prevents an inference
that every conjugated object is absent from the family. The
proved obstruction is internal to one fixed labelled module.

There are already self-adjoint idempotents in the theta data.
Their specified finite group is a product of
$\mathrm{GL}_r(\mathbb F_{g_i})$ over the two auxiliary level
primes and their places, and the central idempotents select
cuspidal finite-level types. This group and action are distinct
from the proposed deck group $\Gal(F/F_0)$. The inspected
definitions supply no identification of them.

The column map is assembled from a theta unfolding map,
determinant characters, and products indexed by the embeddings,
with exponents fixed by the Harish--Chandra parameters. A deck
construction would have to transport these ingredients, the
Hermitian forms, class-lattice weights, integral lattices, and
descended theta and period maps coherently. Uniform valuation
inclusions would then pass as in (10.8), provided the projected
source is the one on which they hold. The exact prime-power
trace identity and its normalization would still have to be
identified. They do not follow from positivity alone.

Finally the literal centering term is also specific: its
\texttt{chiK} is the quadratic Dirichlet character modulo $d$,
and its period term uses
\[
 \log\Omega=\gamma+\log(\pi/v)
                     -\frac{L'(1,\chi_{-d})}{L(1,\chi_{-d})}.
                                                        \tag{11.5}
\]
Replacing this by a statistic of $J_0$ requires a new equality
of the local gain, archimedean terms, and centering constant.
The defining CM quadratic field of this literal construction
is $\Q(\sqrt{-d})$, not $K$. A construction over $K$ is a
modified theorem with its own hypotheses. The declaration that
all roots of unity are $\pm1$ is only an optional proposition in
the inspected host definitions; it is not used as a premise or
an obstruction in this paper.
